% phys.tex -- two colliding discs: from Newton's law to the impulse used by
% imp2d (see ../CONCEPTS.md).  Build:  pdflatex phys.tex  (twice for refs)
\documentclass[a4paper,11pt]{article}

\usepackage[margin=2.3cm]{geometry}
\usepackage{amsmath,amssymb,bm}
\usepackage{booktabs}
\usepackage{enumitem}
\usepackage{tikz}
\usetikzlibrary{arrows.meta,calc}

\setlist{itemsep=2pt,topsep=4pt}

\newcommand{\x}{\bm{x}}

\title{\bfseries Euler rotation equation}
\author{}
\date{}

\begin{document}
\maketitle
\vspace{-2.6cm}

\section{Angular momentum theorem from Newton’s laws}

It follows from Newton’s second law applied to every particle of a rigid body,
plus the rigid-body velocity constraint. The key point is that \(I\) is
constant only in the \emph{body-fixed frame}, so when you take the inertial time
derivative of \(\mathbf L=I\boldsymbol\omega\), a rotating-frame term appears.

For each particle $i$,

\[
m_i \ddot{\mathbf r}_i=\mathbf F_i \qquad \mathbf r_i=\mathbf R_i-\mathbf O.
\]



$\mathbf R_i$ --- position vector of particle $i$ measured from some fixed
reference origin (often an inertial-frame origin).  Let be $O = (0, 0)$ for simplicity.

Cross with \(\mathbf r_i\) and sum over all particles:

\begin{equation}
\sum_i \mathbf r_i\times m_i\ddot{\mathbf r}_i
=
\sum_i \mathbf r_i\times \mathbf F_i .
\label{eq:newton_cross}
\end{equation}

Left-hand side: define the total angular momentum about the chosen origin $O$ as

\[
\mathbf L
=
\sum_i \mathbf r_i\times \mathbf p_i
=
\sum_i \mathbf r_i\times m_i\dot{\mathbf r}_i,
\qquad
\mathbf p_i
=
m_i\dot{\mathbf r}_i.
\]

Now differentiate \(\mathbf L\) with respect to time. Since \(m_i\) is constant,

\[
\frac{d\mathbf L}{dt}
=
\sum_i \frac{d}{dt}\left(\mathbf r_i\times m_i\dot{\mathbf r}_i\right).
\]

Use the product rule for a cross product:

\[
\left[
\quad
\frac{d}{dt}(\mathbf a\cdot\mathbf b)
=
\dot{\mathbf a}\cdot\mathbf b+\mathbf a\cdot\dot{\mathbf b},
\qquad
\frac{d}{dt}(\mathbf a\times\mathbf b)
=
\dot{\mathbf a}\times\mathbf b+\mathbf a\times\dot{\mathbf b}
\quad
\right],
\]

\[
\frac{d}{dt}(\mathbf r_i\times m_i\dot{\mathbf r}_i)
=
\dot{\mathbf r}_i\times m_i\dot{\mathbf r}_i
+
\mathbf r_i\times m_i\ddot{\mathbf r}_i.
\]

The first term vanishes because a vector crossed with itself is
zero (\(\mathbf a\times\mathbf a=0\)):

\[
\dot{\mathbf r}_i\times m_i\dot{\mathbf r}_i
=
m_i(\dot{\mathbf r}_i\times \dot{\mathbf r}_i)
=
0.
\]

Therefore, LHS of \eqref{eq:newton_cross} is

\[
\frac{d\mathbf L}{dt}
=
\sum_i \mathbf r_i\times m_i\ddot{\mathbf r}_i.
\]

The right-hand side of \eqref{eq:newton_cross} is the sum of moments of
all forces about the chosen origin (torque):

\[
\boldsymbol\tau
=
\sum_i \mathbf r_i\times \mathbf F_i.
\]

Now split forces into external and internal:

\[
\boldsymbol\tau
=
\sum_i \mathbf r_i \times \mathbf F_i
=
\sum_i \mathbf r_i \times \mathbf F_i^{\rm ext}
+
\sum_i \mathbf r_i \times \mathbf F_i^{\rm int}
=
\boldsymbol\tau_{\rm ext} + \boldsymbol\tau_{\rm int}
\]

For an ideal rigid body, the internal-force term cancels in pairs.  Standard
assumption: internal forces are equal and opposite along the line joining the
particles:

\[
  \boldsymbol\tau_{\mathrm{int}}
  = \sum_i \mathbf r_i \times \mathbf F_i^{\mathrm{int}}
  = \sum_{i<j} (\mathbf r_i-\mathbf r_j)\times \mathbf F_{ij}
  = \mathbf 0,
\]

because internal forces are equal and opposite and act along the line
joining the particles (strong form of Newton's third law).  So

\[
\boldsymbol\tau
=
\boldsymbol\tau_{\rm ext} + \boldsymbol\tau_{\rm int}
=
\boldsymbol\tau_{\rm ext} + \mathbf 0
=
\boldsymbol\tau_{\rm ext}.
\]

Thus the step gives the angular momentum theorem:

\begin{equation}
\frac{d\mathbf L}{dt}\bigg|_{\text{inertial}}
=
\boldsymbol\tau.
\label{eq:angular_momentum}
\end{equation}

This is often called Euler’s second law, but it is derived from Newton’s laws.

\section{Rigid-body angular momentum: \(\mathbf L=I\boldsymbol\omega\)}

Let the centre of mass be

\[
\mathbf R=\frac{1}{M}\sum_i m_i\mathbf r_i,
\qquad M=\sum_i m_i,
\]

and define the position of particle \(i\) relative to the centre of mass as

\[
\boldsymbol\rho_i=\mathbf r_i-\mathbf R.
\]

Then

\[
\mathbf r_i=\mathbf R+\boldsymbol\rho_i,
\qquad
\dot{\mathbf r}_i=\dot{\mathbf R}+\dot{\boldsymbol\rho}_i.
\]

The angular momentum about the origin (by definition) is

\[
\mathbf L
=
\sum_i \mathbf r_i\times m_i\dot{\mathbf r}_i.
\]

 Substitute into the definition

\[
\mathbf L
=
\sum_i \mathbf r_i\times m_i\dot{\mathbf r}_i
=
\sum_i (\mathbf R+\boldsymbol\rho_i)\times m_i(\dot{\mathbf R}+\dot{\boldsymbol\rho}_i).
\]

Expand the cross product (distribute): for each $i$,

\[
(\mathbf R + \boldsymbol\rho_i) \times (\dot{\mathbf R} + \dot{\boldsymbol\rho}_i)
=
\mathbf R \times (\dot{\mathbf R} + \dot{\boldsymbol\rho}_i)
+
\boldsymbol\rho_i \times (\dot{\mathbf R} + \dot{\boldsymbol\rho}_i).
\]

\[
=
\mathbf R\times\dot{\mathbf R}
+\mathbf R\times\dot{\boldsymbol\rho}_i
+\boldsymbol\rho_i\times\dot{\mathbf R}
+\boldsymbol\rho_i\times\dot{\boldsymbol\rho}_i.
\]

Multiply by $m_i$ and sum over $i$:

\[
\mathbf L
=
\underbrace{\sum_i m_i\,\mathbf R\times\dot{\mathbf R}}_{(1)}
+\underbrace{\sum_i m_i\,\mathbf R\times\dot{\boldsymbol\rho}_i}_{(2)}
+\underbrace{\sum_i m_i\,\boldsymbol\rho_i\times\dot{\mathbf R}}_{(3)}
+\underbrace{\sum_i m_i\,\boldsymbol\rho_i\times\dot{\boldsymbol\rho}_i}_{(4)}.
\]

Evaluate each term.

\begin{enumerate}
\item $\mathbf R$ and $\dot{\mathbf R}$ do not depend on \(i\), so they pull out of the sum:
\[
\sum_i m_i\,\mathbf R\times\dot{\mathbf R}
=
\mathbf R\times\dot{\mathbf R}\sum_i m_i
=
M\,\mathbf R\times\dot{\mathbf R}.
\]

\item $\mathbf R$ is independent of $i$, so

\[
\sum_i m_i\,\mathbf R\times\dot{\boldsymbol\rho}_i
=
\mathbf R\times\sum_i m_i\dot{\boldsymbol\rho}_i
=
\mathbf R\times\mathbf 0
=
\mathbf 0,
\]

because \(\sum_i m_i\dot{\boldsymbol\rho}_i=\dfrac{d}{dt}\sum_i m_i\boldsymbol\rho_i=\mathbf 0\) (see below).

\item $\dot{\mathbf R}$ is independent of $i$, so

\[
\sum_i m_i\,\boldsymbol\rho_i\times\dot{\mathbf R}
=
\left(\sum_i m_i\boldsymbol\rho_i\right)\times\dot{\mathbf R}
=
\mathbf 0\times\dot{\mathbf R}
=
\mathbf 0.
\]

\item This is the only term that survives untouched:

\[
\sum_i m_i\,\boldsymbol\rho_i\times\dot{\boldsymbol\rho}_i.
\]

\end{enumerate}

Why the two vanishing sums are zero?  By definition of the centre of mass,

\[
\sum_i m_i\boldsymbol\rho_i
=
\sum_i m_i(\mathbf r_i-\mathbf R)
=
\sum_i m_i\mathbf r_i-M\mathbf R
=
M\mathbf R-M\mathbf R
=
\mathbf 0.
\]

Differentiating with respect to time (masses constant):

\[
\sum_i m_i\dot{\boldsymbol\rho}_i
=
\frac{d}{dt}\sum_i m_i\boldsymbol\rho_i
=
\mathbf 0.
\]

So both \(\sum_i m_i\boldsymbol\rho_i\) and \(\sum_i
m_i\dot{\boldsymbol\rho}_i\) vanish identically — this is the mathematical
content of "the centre of mass is the mass-weighted mean position."

Putting the four terms together,

\[
\mathbf L
=
M\,\mathbf R\times\dot{\mathbf R}
+\mathbf 0+\mathbf 0
+\sum_i m_i\,\boldsymbol\rho_i\times\dot{\boldsymbol\rho}_i.
\]

Hence

\[
\boxed{
\mathbf L
=
\underbrace{M\,\mathbf R\times\dot{\mathbf R}}_{\text{orbital (CM about origin)}}
+
\underbrace{\sum_i m_i\,\boldsymbol\rho_i\times\dot{\boldsymbol\rho}_i}_{\mathbf L_{\rm CM}\ \text{(spin about CM)}}
}
\]

and therefore

\[
\mathbf L_{\rm CM}
=
\sum_i m_i\,\boldsymbol\rho_i\times\dot{\boldsymbol\rho}_i
=
\mathbf L_{\rm origin}-\mathbf R\times\mathbf P,
\qquad
\mathbf P=M\dot{\mathbf R}.
\]

That is the "substitute the decompositions" step written out in full:
distribute the cross product, pull origin-independent factors out of the sum,
and kill the two cross terms using \(\sum_i m_i\boldsymbol\rho_i=\sum_i
m_i\dot{\boldsymbol\rho}_i=\mathbf 0\).


\hrulefill

But by the definition of the centre of mass,

\[
\sum_i m_i\boldsymbol\rho_i=0,
\qquad
\sum_i m_i\dot{\boldsymbol\rho}_i=0.
\]

So the cross terms vanish, leaving

\[
\mathbf L_{\rm origin}
=
M\mathbf R\times\dot{\mathbf R}
+
\sum_i m_i\,\boldsymbol\rho_i\times\dot{\boldsymbol\rho}_i.
\]

The first term is the angular momentum of the centre of mass about the origin. The second term is the angular momentum **about the centre of mass**:

\[
\boxed{
\mathbf L_{\rm CM}
=
\sum_i m_i\,\boldsymbol\rho_i\times\dot{\boldsymbol\rho}_i
}
\]

Equivalently,

\[
\mathbf L_{\rm CM}
=
\mathbf L_{\rm origin}-\mathbf R\times \mathbf P,
\qquad
\mathbf P=M\dot{\mathbf R}.
\]

So the formula

\[
\mathbf L=\sum_i m_i\,\boldsymbol\rho_i\times\dot{\boldsymbol\rho}_i
\]

is just the angular momentum about the centre of mass, written using coordinates measured from the centre of mass. In that frame, for a rigid body,

\[
\dot{\boldsymbol\rho}_i=\boldsymbol\omega\times\boldsymbol\rho_i,
\]

which gives

\[
\mathbf L
=
\sum_i m_i\,\boldsymbol\rho_i\times(\boldsymbol\omega\times\boldsymbol\rho_i)
=
I\boldsymbol\omega.
\]

\hrulefill

About the center of mass, write a position of particle \(i\) relative to
the centre of mass

\[
\boldsymbol\rho_i=\mathbf r_i-\mathbf R,
\]

where \(\mathbf R\) is the center of mass. For a rigid body,

\[
\dot{\boldsymbol\rho}_i=\boldsymbol\omega\times\boldsymbol\rho_i.
\]

Therefore

\[
\mathbf L
=
\sum_i m_i \boldsymbol\rho_i\times\dot{\boldsymbol\rho}_i
=
\sum_i m_i \boldsymbol\rho_i\times(\boldsymbol\omega\times\boldsymbol\rho_i).
\]

Using

\[
\mathbf a\times(\mathbf b\times\mathbf c)
=
\mathbf b(\mathbf a\cdot\mathbf c)-\mathbf c(\mathbf a\cdot\mathbf b),
\]

we get

\[
\mathbf L
=
\left[
\sum_i m_i\left(|\boldsymbol\rho_i|^2\mathbf 1
-\boldsymbol\rho_i\boldsymbol\rho_i^T\right)
\right]\boldsymbol\omega.
\]

Define the inertia tensor

\[
I=\sum_i m_i\left(|\boldsymbol\rho_i|^2\mathbf 1
-\boldsymbol\rho_i\boldsymbol\rho_i^T\right).
\]

Then

\[
\mathbf L=I\boldsymbol\omega.
\]

In the \emph{body-fixed frame}, the particle coordinates \(\boldsymbol\rho_i\) are
constant, so \(I\) is a constant matrix. In the inertial frame, \(I(t)\) would
generally change with time.

\section{Differentiate \(\mathbf L=I\boldsymbol\omega\) correctly}

The torque equation is inertial:

\[
\boldsymbol\tau=
\frac{d\mathbf L}{dt}\bigg|_{\text{inertial}}.
\]

For any vector \(\mathbf A\), the derivative in the inertial frame is related
to the derivative in a frame rotating with angular velocity
\(\boldsymbol\omega\) by

\[
\frac{d\mathbf A}{dt}\bigg|_{\text{inertial}}
=
\frac{d\mathbf A}{dt}\bigg|_{\text{body}}
+
\boldsymbol\omega\times\mathbf A.
\]

Apply this to \(\mathbf L=I\boldsymbol\omega\). Since \(I\) is constant in the body frame,

\[
\frac{d\mathbf L}{dt}\bigg|_{\text{body}}
=
I\dot{\boldsymbol\omega}.
\]

Therefore

\[
\frac{d\mathbf L}{dt}\bigg|_{\text{inertial}}
=
I\dot{\boldsymbol\omega}
+
\boldsymbol\omega\times(I\boldsymbol\omega).
\]

Using \(\boldsymbol\tau=d\mathbf L/dt|_{\text{inertial}}\),

\[
\boxed{
\boldsymbol\tau
=
I\dot{\boldsymbol\omega}
+
\boldsymbol\omega\times(I\boldsymbol\omega)
}
\]

This is the Euler rotation equation.

\section{Principal-axis form}

If the body axes are principal axes, \(I=\operatorname{diag}(I_1,I_2,I_3)\),
and the vector equation becomes the three Euler equations:

\[
\begin{aligned}
\tau_1 &= I_1\dot\omega_1+(I_3-I_2)\omega_2\omega_3,\\
\tau_2 &= I_2\dot\omega_2+(I_1-I_3)\omega_3\omega_1,\\
\tau_3 &= I_3\dot\omega_3+(I_2-I_1)\omega_1\omega_2.
\end{aligned}
\]

So the extra \(\boldsymbol\omega\times(I\boldsymbol\omega)\) term is not an
additional physical torque; it is the correction that appears because the
body-fixed basis rotates with the body.

\hrulefill

Table of variables used in the derivation

\begin{description}
  \item[\(t\)] time; independent variable for all time derivatives (s).
  \item[\(i,j\)] particle indices labelling individual point masses.
  \item[\(m_i\)] mass of particle \(i\) (kg).
  \item[\(\mathbf r_i\)] position vector of particle \(i\) relative to the chosen origin (m).
  \item[\(\dot{\mathbf r}_i\)] velocity of particle \(i\), \(\dot{\mathbf r}_i = d\mathbf r_i/dt\) (m/s).
  \item[\(\ddot{\mathbf r}_i\)] acceleration of particle \(i\), \(\ddot{\mathbf r}_i = d^2\mathbf r_i/dt^2\) (m/s\(^2\)).
  \item[\(\mathbf R\)] centre-of-mass position of the rigid body (m).
  \item[\(\boldsymbol\rho_i\)] position of particle \(i\) relative to the centre of mass, \(\boldsymbol\rho_i = \mathbf r_i - \mathbf R\) (m).
  \item[\(\mathbf p_i\)] linear momentum of particle \(i\), \(\mathbf p_i = m_i\dot{\mathbf r}_i\) (kg\(\cdot\)m/s).
  \item[\(\mathbf F_i\)] total force acting on particle \(i\) (N).
  \item[\(\mathbf F_i^{\mathrm{ext}}\)] external force acting on particle \(i\) (N).
  \item[\(\mathbf F_{ij}\)] internal force on particle \(i\) due to particle \(j\) (N).
  \item[\(\boldsymbol\tau\)] net external torque about the chosen origin (N\(\cdot\)m).
  \item[\(\mathbf L\)] angular momentum about the chosen origin, \(\mathbf L = \sum_i \mathbf r_i \times m_i\dot{\mathbf r}_i\) (kg\(\cdot\)m\(^2\)/s).
  \item[\(\boldsymbol\omega\)] angular velocity vector of the body-fixed frame (rad/s).
  \item[\(\dot{\boldsymbol\omega}\)] angular acceleration vector (rad/s\(^2\)).
  \item[\(\omega_1,\omega_2,\omega_3\)] components of \(\boldsymbol\omega\) along the principal body axes (rad/s).
  \item[\(\tau_1,\tau_2,\tau_3\)] components of \(\boldsymbol\tau\) along the principal body axes (N\(\cdot\)m).
  \item[\(I\)] inertia tensor, a \(3\times3\) symmetric matrix mapping \(\boldsymbol\omega \mapsto \mathbf L\) (kg\(\cdot\)m\(^2\)).
  \item[\(I_{xx},I_{yy},I_{zz}\)] moments of inertia (diagonal entries), e.g.\ \(I_{xx}=\int\rho(y^2+z^2)\,dV\) (kg\(\cdot\)m\(^2\)).
  \item[\(I_{xy},I_{xz},I_{yz}\)] products of inertia (off-diagonal entries), e.g.\ \(I_{xy}=-\int\rho xy\,dV\) (kg\(\cdot\)m\(^2\)).
  \item[\(I_1,I_2,I_3\)] principal moments of inertia; eigenvalues of \(I\), diagonal in the principal-axis frame (kg\(\cdot\)m\(^2\)).
  \item[\(\mathbf n\)] unit vector along an arbitrary rotation axis.
  \item[\(I_{\mathbf n}\)] moment of inertia about the axis \(\mathbf n\), \(I_{\mathbf n}=\mathbf n^{T}I\,\mathbf n\) (kg\(\cdot\)m\(^2\)).
  \item[\(\mathbf 1\)] \(3\times3\) identity matrix, used in the definition of \(I\).
  \item[\(\rho(\mathbf r)\)] mass density at position \(\mathbf r\) (kg/m\(^3\)).
  \item[\(dV\)] infinitesimal volume element (m\(^3\)).
  \item[\(x,y,z\)] Cartesian coordinates of \(\mathbf r\) (m).
  \item[\(T_{\mathrm{rot}}\)] rotational kinetic energy, \(T_{\mathrm{rot}}=\tfrac12\boldsymbol\omega\cdot I\boldsymbol\omega\) (J).
  \item[\(\mathbf A\)] generic vector used in the rotating-frame derivative rule.
\end{description}

\end{document}
